% !TEX program = pdflatex
\documentclass[11pt]{article}
\usepackage[a4paper,margin=1in]{geometry}
\usepackage{amsmath,amsthm,amssymb,mathtools}
\usepackage{mathrsfs}
\usepackage{bm}
\usepackage{hyperref}
\usepackage{enumitem}
\usepackage{tcolorbox}
\tcbset{colback=white,colframe=black!15}

% Theorem styles
\newtheorem{theorem}{Theorem}[section]
\newtheorem{proposition}[theorem]{Proposition}
\newtheorem{lemma}[theorem]{Lemma}
\newtheorem{corollary}[theorem]{Corollary}
\theoremstyle{definition}
\newtheorem{definition}[theorem]{Definition}
\newtheorem{example}[theorem]{Example}
\theoremstyle{remark}
\newtheorem{remark}[theorem]{Remark}

% Shortcuts
\newcommand{\N}{\mathbb{N}}
\newcommand{\Z}{\mathbb{Z}}
\newcommand{\Q}{\mathbb{Q}}
\newcommand{\R}{\mathbb{R}}
\newcommand{\C}{\mathbb{C}}
\newcommand{\PP}{\mathbb{P}} % primes
\newcommand{\eps}{\varepsilon}
\newcommand{\1}{\mathbf{1}}
\newcommand{\ol}[1]{\overline{#1}}

% Status box
\newenvironment{statusbox}{%
  \begin{tcolorbox}
  \textbf{Status (mainstream vs. USO)}\\[-0.5em]
  \begin{itemize}[leftmargin=1.2em,itemsep=0.25em]
}{\end{itemize}\end{tcolorbox}}

% Title
\title{A Reformulation of  \\ The Unified Spectral Operator (USO) Program\\\large Encodings for $\Pi^0_2$ Statements and Knot Invariants with an Unconditional Result}
\author{(Draft for community review)}
\date{\today}

\begin{document}
\maketitle

\begin{abstract}
We propose a restrained and mathematically explicit version of the \emph{Unified Spectral Operator (USO)} program. The goal is not to solve major conjectures, but to present a concrete operator-theoretic \emph{framework} that encodes (i) certain $\Pi^0_2$ statements and (ii) selected knot invariants into spectral data of bounded linear operators on separable Hilbert spaces. The primary contributions are: (1) explicit definitions of the encoding maps $\Psi$ (for logical statements) and $\Phi$ (for knots); (2) construction principles for the associated operators; and (3) an unconditional proposition demonstrating nontrivial spectral separation for a simple cross-domain encoding. A small ``Collatz-lite'' case study illustrates the pipeline \emph{Definitions $\to$ Operator setup $\to$ Spectral functional $\to$ Claimed implication $\to$ Caveats}. Speculative ideas are quarantined in an appendix and carry no mathematical claims.
\end{abstract}

\begin{tcolorbox}
\textbf{Scope and non-claims.} This paper presents a \emph{framework}, not a suite of proofs. We explicitly do not claim to resolve the Riemann Hypothesis, the Hodge Conjecture, the Yang--Mills mass gap problem, or any Millennium Prize Problem. Our aims are: to define concrete encodings $\Psi$ and $\Phi$, to show how associated spectral operators are constructed, and to prove at least one unconditional proposition evidencing viability. Speculative connections outside pure mathematics are contained only in Appendix~A.
\end{tcolorbox}

\section{Introduction}
A recurring theme across analytic number theory, dynamics, and low-dimensional topology is the translation of discrete structure into spectral information of linear (or transfer) operators. The USO program aspires to provide a unified, computable lens for such translations, while avoiding overreach. We prioritize explicit constructions and modest claims.

\paragraph{Programmatic goals.} Given a class of statements $\mathcal{S}$ (e.g., $\Pi^0_2$) and a class of topological objects $\mathcal{K}$ (e.g., knots), we seek encodings
\[
  \Psi: \mathcal{S} \to \mathsf{Op}(H),\qquad \Phi: \mathcal{K} \to H,
\]
into a fixed separable Hilbert space $H$ and a class of bounded operators $\mathsf{Op}(H)$ so that truth, structure, or invariants correspond to spectral properties (norms, spectra, resolvents, or spectral measures). The emphasis is on \emph{what can actually be proved}.

\begin{statusbox}
  \item \emph{Mainstream status:} RH, Hodge, YM mass gap, Goldbach, Twin Primes, etc. remain \textbf{open}.
  \item \emph{USO claim status:} Proposed encodings + one unconditional proposition (this paper). No claims of problem resolution.
  \item \emph{Speculative content:} Isolated in Appendix~A with disclaimers.
\end{statusbox}

\subsection*{Pipeline}
We standardize the exposition via: \emph{Definitions $\to$ Operator setup $\to$ Spectral functional $\to$ Claimed implication $\to$ Caveats}. Each case study adheres to this structure.

\section{Framework and notation}
Let $H$ be a fixed separable Hilbert space. In concrete instantiations we will take $H=\ell^2(I)$ with a countable index set $I$ (e.g., $I=\N$ or the set $\PP$ of primes), equipped with the standard inner product. Write $\mathsf{B}(H)$ for bounded linear operators on $H$.

\begin{definition}[Encoding maps]
\label{def:encodings}
\begin{enumerate}[label=(\alph*),leftmargin=2em]
  \item ($\Pi^0_2$ encoding) An encoding is a map $\Psi$ from a statement $S\in\mathcal{S}$ to a family of operators $\{U_{\sigma}(S)\}_{\sigma\in\Sigma}\subset \mathsf{B}(H)$ together with a specified spectral predicate $\mathcal{P}$ (e.g., ``$\rho(U_{\sigma}(S))\le 1$'' or ``$\exists$ fixed point'') such that $S$ \emph{implies or is implied by} $\mathcal{P}(U_{\sigma}(S))$.
  \item (Knot encoding) An encoding is a map $\Phi: \mathcal{K}\to H$ taking a knot $K$ to a concrete vector of invariants in the chosen basis. Typical choices include: coefficients of the Jones or HOMFLY--PT polynomial in a fixed monomial basis; values of a finite set of normalized invariants; or data from the $\mathrm{SL}_2(\C)$ character variety embedded in a finite coordinate chart.
\end{enumerate}
\end{definition}

\begin{remark}
The \emph{USO} perspective is not a single universal operator but a \emph{template}: in each application one explicit operator (or a small family) is constructed with clearly stated domain/codomain and properties. We avoid bundled opaque hypotheses.
\end{remark}

\section{Operator setup: principles and properties}
For any operator introduced, we verify linearity, boundedness, and (when used) self-adjointness or compactness. We avoid assuming compact self-adjointness without proof.

\begin{definition}[Spectral functionals]
Given $T\in\mathsf{B}(H)$, typical spectral functionals used are: the operator norm $\|T\|$, the spectral radius $\rho(T)$, point spectrum $\sigma_p(T)$, and the full spectrum $\sigma(T)$. When $T$ is normal, $\|T\|=\rho(T)$. We also consider compression norms $\|PTP\|$ for orthogonal projections $P$.
\end{definition}

\section{Unconditional result: a simple cross-domain spectral separation}
Let $\PP$ denote the set of primes and set $H=\ell^2(\PP)$. For a hyperbolic knot $K$ with positive hyperbolic volume $\mathrm{vol}(K)>0$, define a vector $\Phi_{\mathrm{vol}}(K)\in H$ by
\[
  \Phi_{\mathrm{vol}}(K)_p := p^{-\mathrm{vol}(K)} \qquad (p\in\PP).
\]
Define the associated diagonal operator $D_{\Phi_{\mathrm{vol}}(K)}\in\mathsf{B}(H)$ by $D_{\Phi_{\mathrm{vol}}(K)} e_p = \Phi_{\mathrm{vol}}(K)_p\, e_p$, where $\{e_p\}_{p\in\PP}$ is the standard orthonormal basis.

\begin{proposition}[Spectral distinction via prime-indexed diagonal encodings]
\label{prop:diagonal}
Let $K_1,K_2$ be hyperbolic knots with $\mathrm{vol}(K_1)\neq \mathrm{vol}(K_2)$. Then each $D_{\Phi_{\mathrm{vol}}(K_i)}$ is compact, and
\[
  \sigma\big(D_{\Phi_{\mathrm{vol}}(K_i)}\big)=\{0\}\cup\{p^{-\mathrm{vol}(K_i)}:p\in\PP\}.
\]
Moreover,
\[
  \big\|D_{\Phi_{\mathrm{vol}}(K_1)}-D_{\Phi_{\mathrm{vol}}(K_2)}\big\|=\sup_{p\in\PP}\big|p^{-\mathrm{vol}(K_1)}-p^{-\mathrm{vol}(K_2)}\big|=\big|2^{-\mathrm{vol}(K_1)}-2^{-\mathrm{vol}(K_2)}\big|>0,
\]
so the spectra differ beyond the shared point $\{0\}$.
\end{proposition}
\begin{proof}
Since $\mathrm{vol}(K_i)>0$, we have $\Phi_{\mathrm{vol}}(K_i)_p\to 0$ as $p\to\infty$. A diagonal operator on $\ell^2$ is compact iff its diagonal entries tend to zero, hence compactness. The spectrum of a diagonal operator on $\ell^2$ is the closure of its set of diagonal entries, which here is exactly $\{0\}\cup\{p^{-\mathrm{vol}(K_i)}:p\in\PP\}$. For the norm of the difference of diagonal operators we have $\|\mathrm{diag}(a_p)-\mathrm{diag}(b_p)\|=\sup_p|a_p-b_p|$. The function $p\mapsto p^{-v}$ is strictly decreasing in $p$ for fixed $v>0$, so the supremum over primes is attained at $p=2$. If $\mathrm{vol}(K_1)\neq\mathrm{vol}(K_2)$ the value is strictly positive, and the spectral sets differ (beyond $\{0\}$).
\end{proof}

\begin{remark}
Proposition~\ref{prop:diagonal} is intentionally simple: it shows a nontrivial, unconditional spectral separation obtained from a concrete cross-domain encoding (knot volume $\to$ prime-indexed diagonal). Stronger encodings may replace volume by a vetted invariant vector and diagonal by bounded Toeplitz or integral operators; compactness and spectral description then require separate proofs.
\end{remark}
\begin{proposition}[Toeplitz-compression invariant of a knot-invariant subspace]\label{prop:toeplitz}
Let \(\Phi\) map each knot \(K\) to the vector of coefficients of its Jones polynomial in the fixed monomial basis, and let \(V=\overline{\mathrm{span}}\{\Phi(K):K\}\subset \ell^2(\Z)\).
If \(T\) is a bounded self-adjoint operator on \(\ell^2(\Z)\) that commutes with the shift, then the compression norm \(\|P_V T P_V\|\) is a well-defined invariant of the subspace \(V\), independent of the ambient completion.
\end{proposition}
\begin{proof}[Proof sketch]
Commutation with the shift makes \(T\) a (self-adjoint) Toeplitz operator. Closed shift-invariant subspaces are invariant under Toeplitz operators, and compressions to such subspaces preserve boundedness and self-adjointness. Hence \(\|P_V T P_V\|\) depends only on \(V\).
\end{proof}

\section{Toy case study: a Collatz-lite encoding}
We now present a concrete transfer-style operator \(L_s\) (with subscript \(s\) exactly as written) on \(H=\ell^2(\N)\) that captures the weighted preimage structure of the Collatz map. No conjectural equivalences are asserted.

\subsection*{Definitions}
Let \(C:\N\to\N\) be the usual Collatz map
\[
  C(n)=\begin{cases}
    n/2, & n\text{ even},\\
    3n+1, & n\text{ odd}.
  \end{cases}
\]
For parameters \(\alpha,\beta\in[0,1)\) and real \(s\ge 0\), define the linear operator
\[
  L_s: \ell^2(\N)\longrightarrow \ell^2(\N),\qquad (L_s f)(n) := \alpha^{\,s} f(2n)\;+\;\beta^{\,s}\, \mathbf{1}_{\{n\equiv 1\ (3)\}}\, f\!\left(\tfrac{n-1}{3}\right).
\]
Equivalently, \(L_s\) sends each value \(f(m)\) forward to its image index under \(C\), with weight \(\alpha^{s}\) for even-preimage transitions (\(m=2n\)) and \(\beta^{s}\) for odd-preimage transitions (\(m=(n-1)/3\) when \(n\equiv1\pmod 3\)).

\subsection*{Boundedness}
\begin{lemma}[Boundedness of \(L_s\)]\label{lem:Ls-bounded}
For any \(s\ge 0\) and \(\alpha,\beta\in[0,1)\), the operator \(L_s\in \mathsf{B}(\ell^2(\N))\) and
\[
  \|L_s\|\;\le\; \alpha^{\,s} + \beta^{\,s}.
\]
\end{lemma}
\begin{proof}
Let \((S_2 f)(n):=f(2n)\) and \((T_3 f)(n):=\mathbf{1}_{\{n\equiv1\ (3)\}} f((n-1)/3)\). Then \(L_s=\alpha^{s} S_2 + \beta^{s} T_3\). We have
\(\sum_{n\ge1}|(S_2 f)(n)|^2=\sum_{n\ge1}|f(2n)|^2\le\sum_{m\ge1}|f(m)|^2\), hence \(\|S_2\|\le1\).
Similarly, the map \(n\mapsto (n-1)/3\) on \(\{n\equiv1\ (3)\}\) is injective into \(\N\), so
\(\sum_{n\ge1}|(T_3 f)(n)|^2\le\sum_{m\ge1}|f(m)|^2\) and \(\|T_3\|\le1\).
Therefore \(\|L_s\|\le \alpha^{s}\|S_2\|+\beta^{s}\|T_3\|\le \alpha^{s}+\beta^{s}\).
\end{proof}

\subsection*{What the spectrum measures}
The operator \(L_s\) is a \emph{weighted preimage transfer}: its action aggregates contributions from immediate preimages under \(C\).
The spectral radius \(\rho(L_s)\) upper-bounds the exponential growth rate of weighted preimage counts along the Collatz graph
(with weights \(\alpha^{s}\) for even edges and \(\beta^{s}\) for odd edges).
Point spectrum, when present, corresponds to \emph{eigen-distributions} that are stable under the two local moves
\(n\mapsto 2n\) and (when allowed) \(n\mapsto (n-1)/3\).
These interpretations are purely structural and do not assert any equivalence to global reachability properties.

\section*{Proposition X: Spectral Distinction of Knots via Prime Statistics}
\begin{proposition}[Proposition X (Spectral Distinction of Knots via Prime Statistics)]\label{prop:X}
Let \(K_1, K_2\) be two distinct hyperbolic knots with different hyperbolic volumes, \(\mathrm{vol}(K_1) \neq \mathrm{vol}(K_2)\). Let \(\Phi_{\mathrm{vol}}(K)\) be the vector in \(\ell^2(\PP)\) (indexed by primes) defined by \(\Phi_{\mathrm{vol}}(K)_p = p^{-\mathrm{vol}(K)}\). Then the associated diagonal operators \(D_{\Phi_{\mathrm{vol}}(K_1)}\) and \(D_{\Phi_{\mathrm{vol}}(K_2)}\) on \(\ell^2(\PP)\) are compact, and their spectra are disjoint. In particular,
\[
\big\|D_{\Phi_{\mathrm{vol}}(K_1)} - D_{\Phi_{\mathrm{vol}}(K_2)}\big\| \ge \big|p^{-\mathrm{vol}(K_1)} - p^{-\mathrm{vol}(K_2)}\big| \quad \text{for any prime } p.
\]
\end{proposition}
\begin{proof}[Proof sketch]
Diagonal entries \(p^{-\mathrm{vol}(K)}\to 0\) as \(p\to\infty\), hence compactness. The spectrum of a diagonal operator on \(\ell^2\) is the closure of its diagonal entries, so differing volumes produce distinct nonzero spectral values. The operator-norm lower bound follows from \(\|\mathrm{diag}(a_p)-\mathrm{diag}(b_p)\|=\sup_p |a_p-b_p|\); choosing any prime \(p\) yields the stated inequality.
\end{proof}

\section{Related work (positioning)}
The use of operators to encode arithmetic or dynamical information has precedents: transfer (Ruelle--Perron--Frobenius) operators in dynamics; Selberg-type trace formulae; Toeplitz and Hankel operators; and C$^*$-algebraic approaches to zeta-like structures. Our contribution here is a \emph{template} that insists on explicit, computable encodings and on separating unconditional results from heuristic numerics.

\section{Discussion and caveats}
This draft favors clarity over breadth. The encodings chosen are intentionally concrete; stronger choices (and deeper theorems) might follow, but only with precise operator-theoretic control (boundedness, compactness, self-adjointness/normality, approximation schemes, and truncation error bounds).
\section{Integration of the Unified Spectral Operator (USO) with DRMM}
\label{sec:uso-drmm}

\subsection{State Space, Base Operator, and Channels}
Let $H=\ell^2(\mathbb{P})$ with canonical basis $\{e_p\}_{p\in\mathbb{P}}$. 
A USO encoding supplies a bounded, explicitly analyzable base operator 
$X:\!H\to H$. For the seed construction we take a diagonal USO base from a
hyperbolic knot $K_0$ with volume $\mathrm{vol}(K_0)$:
\[
X \,=\, \mathrm{diag}(x_p),\qquad x_p:=p^{-\mathrm{vol}(K_0)}\quad (p\in\mathbb{P}).
\]
We consider a prime–scripted, self-adjoint, parameterized perturbation (the DRMM plant)
\begin{equation}
\label{eq:plant}
U(\omega)\;=\;X+\sum_{p\le p_{\max}} w_p\,B_p(\omega),
\end{equation}
where each channel $B_p(\omega)=B_p(\omega)^\ast$ satisfies certified bounds
\begin{equation}
\label{eq:channel-bounds}
\|B_p(\omega)\|\le b_p,\qquad \|\partial_\omega B_p(\omega)\|\le b'_p,
\end{equation}
with known nonnegative sequences $(b_p)_{p},(b'_p)_{p}$. Denote the weighted
$\ell_1$ seminorms by
\[
\|w\|_{1,b}:=\sum_{p\le p_{\max}} b_p\,|w_p|,\qquad 
\|w\|_{1,b'}:=\sum_{p\le p_{\max}} b'_p\,|w_p|.
\]

\subsection{Lawfulness Budgets (Prime-Gating, Smoothness, Commutators)}
Fix $\alpha>0$ and weights $(\lambda_p)_{p}$. The \emph{prime-gating} budget constrains $w$ to a tube around the lawful profile $p^{-\alpha}$:
\begin{equation}
\label{eq:prime-gating}
R_{\Lambda}^{(\alpha)}(w,\kappa):=\sum_{p\le p_{\max}}\lambda_p\,|w_p-\kappa\,p^{-\alpha}|\ \le\ \tau,
\end{equation}
for some gain $\kappa\in\mathbb{R}$ and tolerance $\tau>0$. A spectral smoothness budget caps the parametric slope:
\begin{equation}
\label{eq:slope-cap}
\|w\|_{1,b'}\ \le\ s_{\max}.
\end{equation}
If channels do not commute, introduce an aggregate commutator budget
\begin{equation}
\label{eq:comm-budget}
\sum_{p<q}\|[B_p(\omega),B_q(\omega)]\|\ \le\ \gamma,
\end{equation}
with $\gamma\ge 0$ (trivial when $B_p$ are diagonal).

\subsection{Certified Gap and Slope}
Let $\{\lambda_k(\omega)\}$ denote the eigenvalues of $U(\omega)$ (counting multiplicity).
For a tracked eigen-cluster of $X$ with minimal separation
\begin{equation}
\label{eq:deltaS}
\delta_S\ :=\ \min_{i\neq j}\,|x_i-x_j|\quad \text{over the tracked indices,}
\end{equation}
define the realized cluster gap of $U(\omega)$ by 
$\mathrm{gap}(U(\omega)):=\min_{i\neq j}|\lambda_i(\omega)-\lambda_j(\omega)|$ over those indices.

\begin{theorem}[Gap--Slope Certificate]
\label{thm:gap-slope}
Suppose $X$ is normal (diagonal above) and $B_p(\omega)$ satisfy \eqref{eq:channel-bounds}.
For any feasible $(w,\kappa)$ under \eqref{eq:prime-gating}--\eqref{eq:comm-budget} and all $\omega$,
\begin{equation}
\label{eq:gap-bound}
\mathrm{gap}(U(\omega))\ \ge\ \delta_S\ -\ 2\,\|w\|_{1,b},
\end{equation}
and for each differentiable eigenvalue branch $\lambda_k(\omega)$ in the cluster,
\begin{equation}
\label{eq:slope-bound}
\bigl|\tfrac{d}{d\omega}\lambda_k(\omega)\bigr|\ \le\ \|w\|_{1,b'}\ \le\ s_{\max}.
\end{equation}
In particular, if $\|w\|_{1,b}<\delta_S/2$ the cluster remains separated for all $\omega$.
\end{theorem}

\begin{proof}[Proof sketch]
Weyl-type perturbation bounds yield $\bigl||\lambda_i-\lambda_j|-|x_i-x_j|\bigr|\le 
\|U-X\|\le \sum_p \|w_p B_p\|\le \|w\|_{1,b}$ twice (upper and lower), giving \eqref{eq:gap-bound}.
For the slope, differentiate \eqref{eq:plant}, apply standard eigenvalue
perturbation estimates and \eqref{eq:channel-bounds} to obtain
$|\tfrac{d}{d\omega}\lambda_k|\le \|\partial_\omega U\|
\le \sum_p |w_p|\,\|\partial_\omega B_p\|
=\|w\|_{1,b'}$, yielding \eqref{eq:slope-bound}.
\end{proof}

\subsection{Convex Controller: A Practical Program}
To widen separation while capping slope, solve for $(w,\kappa)$
\begin{equation}
\label{eq:controller}
\begin{aligned}
\min_{w,\kappa}\quad & 
J(w):= 2\,\|w\|_{1,b}\ +\ \mu\,\|w\|_{1,b'}\\
\text{s.t.}\quad & R_{\Lambda}^{(\alpha)}(w,\kappa)\ \le\ \tau,\\
& \|w\|_{1,b'}\ \le\ s_{\max},\\
& \text{commutator budget \eqref{eq:comm-budget} (data-encoded).}
\end{aligned}
\end{equation}
Here $\mu\ge 0$ trades off nominal gap expansion against slope regularity. 
Given $(b,b',\Lambda,\alpha,\tau,s_{\max},\gamma)$, the solution $(w^\star,\kappa^\star)$ yields
certificates
\[
B_{\mathrm{gap}}:=2\|w^\star\|_{1,b},\qquad B_{\mathrm{slope}}:=\|w^\star\|_{1,b'},
\]
and thus, by Theorem~\ref{thm:gap-slope}, 
$\mathrm{gap}(U)\ge \delta_S-B_{\mathrm{gap}}$ and 
$|d\lambda_k/d\omega|\le B_{\mathrm{slope}}\le s_{\max}$.

\subsection{Proximal Realization (Budget-Preserving)}
\label{subsec:prox}
We employ a simple, auditable proximal scheme.

\paragraph{Inputs} weights $(b,b',\lambda)$, budgets $(\tau,s_{\max},\gamma)$, grid $\alpha\in\mathcal A$.

\paragraph{Iteration} For a fixed $\alpha$:
\begin{enumerate}
\item \textbf{Soft-threshold step:} 
$w \leftarrow \mathrm{SoftThresh}_{\eta(2b+\mu b')}(w - g)$ with any subgradient $g\in\partial J$.
\item \textbf{Budget projection:} Project $(w,\kappa)$ onto the convex set 
$\{(w,\kappa):R_{\Lambda}^{(\alpha)}(w,\kappa)\le \tau,\ \|w\|_{1,b'}\le s_{\max}\}$.
\item \textbf{Fit $\kappa$:} Update $\kappa$ as a weighted median of $w_p / p^{-\alpha}$ under weights $\lambda_p$.
\end{enumerate}
Stop when the primal residual and objective variation fall below tolerance. 
Return $(w^\star,\kappa^\star)$ and $(B_{\mathrm{gap}},B_{\mathrm{slope}})$.

\subsection{Verification and Falsification Protocol}
\label{subsec:verify}
\begin{enumerate}
\item \textbf{Lanczos sweep.} On a grid $\omega\in\Omega$, track the targeted eigenpairs of $U(\omega)$; record realized gap $g_{\mathrm{real}}$ and slope $s_{\mathrm{real}}$.
\item \textbf{Certificate check.} Verify $g_{\mathrm{real}}\ge \delta_S-B_{\mathrm{gap}}$ and $s_{\mathrm{real}}\le B_{\mathrm{slope}}$.
\item \textbf{Null controllers.} Generate $M$ nulls by (i) sign randomization of $w^\star$ or (ii) prime-index shuffles preserving $|w_p^\star|$, while holding budgets equal. 
Compute $g_{\mathrm{real}}^{(m)}$ and report $p\!=\!\frac{1}{M}\#\{m:\ g_{\mathrm{real}}^{(m)}\ge g_{\mathrm{real}}\}$.
\item \textbf{Report.} Publish $(\alpha,\tau,\gamma,s_{\max})$, $(B_{\mathrm{gap}},B_{\mathrm{slope}})$, $(g_{\mathrm{real}},s_{\mathrm{real}},p)$, and logs.
\end{enumerate}

\subsection{Seed Instance (Diagonal, Commuting Channels)}
\label{subsec:seed}
Take $K_0=4_1$ (figure-eight knot; $\mathrm{vol}(K_0)\approx2.03$), so $X=\mathrm{diag}(p^{-\mathrm{vol}(K_0)})$.
Choose a comparator $K_1$ with different volume for secondary spectra.
Let $B_p(\omega)=\beta_p(\omega)E_{pp}$ with $|\beta_p|\le b_p$ and $|\partial_\omega \beta_p|\le b'_p$
(e.g.\ $b_p=b'_p=p^{-1-\varepsilon}$). Then \eqref{eq:comm-budget} is vacuous.
Solve \eqref{eq:controller} on a small grid $\mathcal A=\{0.5,0.7,1.0,1.3\}$ with conservative $(\tau,s_{\max})$.
By Theorem~\ref{thm:gap-slope}, any $(w^\star,\kappa^\star)$ with $\|w^\star\|_{1,b}<\delta_S/2$ certifies separation; 
verification uses a single sweep and null baselines as in \S\ref{subsec:verify}.

\subsection{Reproducibility Checklist}
\begin{itemize}
\item \textbf{Data:} $(x_p)$ from the chosen USO encoding; channel bounds $(b_p,b'_p)$; weights $(\lambda_p)$.
\item \textbf{Budgets:} $(\alpha,\tau,s_{\max},\gamma)$ and truncation $p_{\max}$.
\item \textbf{Outputs:} $(w^\star,\kappa^\star)$, certificates $(B_{\mathrm{gap}},B_{\mathrm{slope}})$, realized $(g_{\mathrm{real}},s_{\mathrm{real}})$, and $p$-value.
\end{itemize}

\begin{statusbox}
  \item \textbf{What is proved here?} Proposition~\ref{prop:diagonal} and Lemma~\ref{lem:finite}; both unconditional.
  \item \textbf{What is not claimed?} Any resolution of major open problems; any equivalence between Collatz and a spectral predicate; any universal compact self-adjoint USO.
  \item \textbf{Where next?} Replace the diagonal example by a vetted invariant vector and nontrivial operator (Toeplitz/compression or integral kernels) and prove compactness and spectral statements; develop a nontruncated Collatz-operator on a natural Banach/Hilbert space.
\end{statusbox}

\section*{Appendix A: Speculative extensions and philosophical context}
\noindent\textit{Disclaimer.} The material in this appendix is \textbf{speculative} and does not support, refine, or imply any theorem in the main text. It records informal analogies that originally motivated the USO program. No claims below are used in proofs, and no terminology here should be read as physical modeling.

\subsection*{A.1. Vocabulary discipline}
Terms such as “topological qubit”, “entanglement”, “Noether flux”, “spectral flux”, and “T3\textsuperscript{-} universality” are metaphors only and are \emph{not} employed in the technical development. In particular:
\begin{itemize}[leftmargin=1.2em]
\item \textbf{No qubits.} We do not define or use any qubit model. Knots are encoded via vectors $\Phi(K)\in H$; these are not quantum states in a physical sense.
\item \textbf{No entanglement.} We do not specify a tensor-factorization $H=H_A\otimes H_B$, hence no notion of entanglement is available or used.
\item \textbf{No Noether claims.} We do not introduce a Lagrangian, a configuration space, or continuous symmetries for a dynamical system. Thus “Noether’s theorem” is inapplicable here.
\end{itemize}

\subsection*{A.2. Why we still note the metaphors}
The operator viewpoint sometimes invites language borrowed from physics (e.g., “symmetry”, “state”), which can help intuition. However, this appendix is strictly a repository for such metaphors to prevent contamination of the main mathematical narrative.

\subsection*{A.3. Non-goals}
We do not propose a model of topological quantum computation, quantum error correction, or quantum dynamics. Any future work in that direction would require: (i) a concrete tensor-factorization, (ii) a gate set and dynamics, (iii) noise/error models, and (iv) rigorous links to established topological quantum computing frameworks. None of these are attempted here.

\subsection*{A.4. Safe rephrasings (if metaphors are mentioned informally)}
If informal language is retained (e.g., in introductions), it must be paired with a mathematical rephrasing. Examples:
\begin{itemize}[leftmargin=1.2em]
\item “SU(2) symmetry preserves entanglement” $\rightsquigarrow$ \emph{Unitary invariance of the inner product on $H$; no tensor-product structure is assumed or used.}
\item “Knots as qubits” $\rightsquigarrow$ \emph{Knots mapped to vectors $\Phi(K)$ in a fixed Hilbert space $H$; linear combinations are formal encodings, not physical superpositions.}
\end{itemize}

\section*{Appendix B: Errata and retractions for \textnormal{USOFisical.pdf}}

\paragraph{B.1 Topological qubits.}
\emph{Retraction.} The phrase “topological qubits” was used without defining a tensor-factorization, gate set, or error-correcting structure. It is removed from the mainline. Knots are encoded only via vectors $\Phi(K)\in H$.

\paragraph{B.2 Entanglement preservation (former Theorem 1).}
\emph{Retraction.} The statement conflated inner-product invariance under unitaries with quantum entanglement, which requires $H=H_A\otimes H_B$. No such factorization is present. The theorem is deleted.
\medskip

\noindent\textbf{Replacement (harmless lemma).}
\begin{lemma}[Unitary invariance of the inner product]
If $U:H\to H$ is unitary, then $\langle Ux,Uy\rangle=\langle x,y\rangle$ for all $x,y\in H$.
\end{lemma}
\begin{proof}
Standard.
\end{proof}

\paragraph{B.3 Noether-Flux (former Theorem 6).}
\emph{Retraction.} Noether’s theorem pertains to continuous symmetries of a Lagrangian/Hamiltonian system on a configuration space; the prior text introduced none. The claim is removed with no replacement.

\paragraph{B.4 Spectral Flux and $T3^{-}$ universality (former Theorem 8).}
\emph{Retraction.} The composite $\Pi=\pi\circ\mathcal{R}_{ABC}\circ P_{inv}\circ J$ was undefined (no $J$, no $\mathcal{R}_{ABC}$). The theorem is deleted.

\paragraph{B.5 Prime-number encoding via $X_P$ (former Theorem 3).}
\emph{Correction.} The matrix defined by $[X_P]_{ij}=\zeta(1/2+i\log(p_i/p_j))$ is a heuristic construction with no proven link to prime distribution or to the spectrum of any canonical operator. All claims of “encoding primes” or “locating zeta zeros in the discrete spectrum” are removed. If retained, $X_P$ is presented only as a toy family with no number-theoretic implications.

\paragraph{B.6 Code quality and honesty.}
\emph{Correction.} Functions returning hard-coded Jones/HOMFLY polynomials or volumes are removed. If computational examples are desired, they must either:
\begin{itemize}[leftmargin=1.2em]
\item implement a bona fide algorithm (e.g., Kauffman bracket for Jones), with explicit complexity/limits; or
\item cite an external database and state clearly that values are looked up, not computed.
\end{itemize}

\paragraph{B.7 Quantum Hilbert space $\mathcal{H}_Q$.}
\emph{Retraction.} The object $\mathcal{H}_Q$ was described as finite-dimensional and spanned by “topological qubits” without specifying a basis or physical semantics. It is removed. Any finite-dimensional subspace used in examples will be denoted by $V\subset H$ with a declared basis and no physics meaning attached.

\section*{Appendix C: Code and data policy (minimal)}
All code and datasets, if included, are strictly exploratory and serve to illustrate definitions already given in theorems/lemmas. No numerical output is framed as proof. For polynomial invariants, either (i) implement a standard algorithm (e.g., Kauffman bracket for Jones) with input diagrams and complexity notes, or (ii) use a curated database with citations, clearly stating values are looked up.


\end{thebibliography}

\end{document}
